% =====================================================================
%  Apêndice D — Declaração de uso de IA generativa
% =====================================================================
\chapter{Declaração de uso de inteligência artificial generativa}
\label{ap:ia}

\begin{table}[H]
\centering
\caption{Uso de IA generativa neste ciclo}
\label{tab:ia}
\footnotesize
\setlength{\tabcolsep}{3.5pt}
\begin{tabular}{@{}L{3.2cm}L{2.6cm}L{5.4cm}L{4.0cm}@{}}
\toprule
\textbf{Etapa} & \textbf{Ferramenta} & \textbf{Para quê} & \textbf{O que a equipe verificou} \\
\midrule
\emph{Firmware} BLE (\emph{commit} \texttt{122c41e}) & Claude (Anthropic) &
  apoio na escrita do código do \emph{firmware} & ensaios com o nRF Connect registrados no README \\
Revisão do repositório & Claude (Anthropic), no Claude Code &
  organizou o repositório para a AP2, movendo o projeto PlatformIO para \caminho{firmware/} &
  o \emph{firmware} em \caminho{firmware/} foi compilado, gravado na placa e está funcionando \\
Gravador e análise de dados & Claude (Anthropic), no Claude Code &
  escreveu \texttt{gravador.py}, \texttt{protocolo.py}, \texttt{qualidade.py}, \texttt{limpar.py}, \texttt{grafico.py}, \texttt{ensaio.py}, \texttt{ct05.py} e os testes automáticos, a partir do contrato BLE do README; corrigiu o gravador durante a campanha (SIGINT e SIGTERM) &
  10 testes automáticos passando; a equipe rodou o gravador com o nó (teste de fumaça, ensaios e campanha de 24/10) \\
Ensaios e campanha & Claude (Anthropic), no Claude Code, no notebook do gravador &
  rodou o teste de fumaça, os ensaios cronometrados (\CT{02}, \CT{13}, \CT{14}), o \CT{05} e os ensaios automáticos \CT{15} e \CT{16}; iniciou, acompanhou e encerrou o gravador; registrou o diário a partir do que a equipe contou &
  as ações físicas foram da equipe: instalação e reposicionamento do nó, botão EN (\CT{17}), cabo USB (\CT{18}) e os eventos do quarto; os dados vêm do nó, sem edição \\
Relatório & Claude (Anthropic), no Claude Code &
  primeira versão do texto deste relatório, a partir do relatório da AP1, do README, do código e das instruções de entrega; depois da campanha, preencheu os números a partir das saídas do \texttt{qualidade.py} e dos arquivos de ensaio &
  a equipe revisou cada capítulo e forneceu os dados que só ela tinha (local, papéis, eventos, objetivo da AP3) \\
Não houve uso em & \multicolumn{3}{L{12.0cm}}{montagem do circuito e coleta dos dados pelo nó.} \\
\bottomrule
\end{tabular}
\fonte{Elaborada pela equipe.}
\end{table}

\vspace{0.4cm}

\noindent\textbf{Declaração.} A equipe declara que o conteúdo deste
relatório e do código a ele associado foi revisado e é compreendido por
todos os integrantes, e que todo uso de ferramentas de inteligência
artificial generativa está registrado na tabela acima. Todos os
integrantes leram esta declaração e confirmam a revisão do relatório.

\vspace{0.4cm}

\imprimirTabelaIntegrantes

\vspace{0.4cm}

\noindent\localEntrega, \dataEntrega.
